% !TEX root = task-6-5-3.tex
% Editable task bank. Append new BankTask blocks at the end of the document.
\documentclass[12pt,a4paper]{article}
\usepackage{physics-textbook}
\usepackage{graphicx,needspace}
\usetikzlibrary{decorations.pathmorphing}
\newif\ifShowSolutions
\ifdefined\TutorEdition
  \ShowSolutionstrue
\else
  \ShowSolutionsfalse % Change to \ShowSolutionstrue to display solutions.
\fi
\newcounter{banktask}
\newsavebox{\BankTaskBox}
\newcommand{\BankTaskContent}[3]{%
  \noindent\textbf{\thebanktask.}\quad #2\par
  \ifShowSolutions
    \Needspace{50mm}% Keep the solution heading with its opening text or diagram.
    #3\par
    \ifdefined\TutorEdition
      \ifcsname TutorSource#1\endcsname
        \nopagebreak\smallskip{\small\csname TutorSource#1\endcsname\par}%
      \fi
    \fi
  \fi
}
\newcommand{\BankTask}[3]{%
  \refstepcounter{banktask}\par\medskip
  \sbox{\BankTaskBox}{\begin{minipage}{\linewidth}%
    \BankTaskContent{#1}{#2}{#3}\end{minipage}}%
  \ifdim\dimexpr\ht\BankTaskBox+\dp\BankTaskBox\relax>\dimexpr\textheight-20mm\relax
    \Needspace{35mm}\BankTaskContent{#1}{#2}{#3}%
  \else
    \Needspace{\dimexpr\ht\BankTaskBox+\dp\BankTaskBox\relax}%
    \noindent\usebox{\BankTaskBox}\par
  \fi
  \medskip
}
\begin{document}
\BookSubsection{6.5.3}{Планетарная модель атома}


% Source sdamgia-7047: https://phys-ege.sdamgia.ru/problem?id=7047
\BankTask{sdamgia-7047}{
Какое утверждение соответствует планетарной модели атома?\par 1) Ядро  — в центре атома, заряд ядра положителен, электроны  — на орбитах вокруг ядра.\par 2) Ядро  — в центре атома, заряд ядра отрицателен, электроны  — на орбитах вокруг ядра.\par 3) Электроны  — в центре атома, ядро обращается вокруг электронов, заряд ядра положителен.\par 4) Электроны  — в центре атома, ядро обращается вокруг электронов, заряд ядра отрицателен.
}{
\smallskip\noindent\textbf{Решение}\par
Планетарной модели атома соответствует модель «ядро  — в центре атома, заряд ядра положителен, электроны  — на орбитах вокруг ядра».

Правильный ответ указан под номером $1$.
}

% Source sdamgia-7079: https://phys-ege.sdamgia.ru/problem?id=7079
\BankTask{sdamgia-7079}{
Какое представление о строении атома соответствует модели атома Резерфорда?\par 1) Ядро  — в центре атома, заряд ядра положителен, бóльшая часть массы атома сосредоточена в электронах.\par 2) Ядро  — в центре атома, заряд ядра отрицателен, бóльшая часть массы атома сосредоточена в электронной оболочке.\par 3) Ядро  — в центре атома, заряд ядра отрицателен, бóльшая часть массы атома сосредоточена в ядре.\par 4) Ядро  — в центре атома, заряд ядра положителен, бóльшая часть массы атома сосредоточена в ядре
}{
\smallskip\noindent\textbf{Решение}\par
Модели Резерфорда соответствует следующая модель: ядро  — в центре атома, заряд ядра положителен, бóльшая часть массы атома сосредоточена в ядре.

Правильный ответ указан под номером $4$.
}

\end{document}
